\documentclass[11pt,a4paper]{article}
\usepackage[utf8]{inputenc}
\usepackage[T1]{fontenc}
\usepackage{amsmath,amssymb}
\usepackage{graphicx}
\usepackage{booktabs}
\usepackage{hyperref}
\usepackage{cite}
\usepackage{geometry}
\geometry{margin=2.5cm}

\title{\textbf{Calibration-Quoted $T_2$ Overestimates Random-Axis Echo $T_2$ on a Drifting Superconducting Transmon: Correlated $T_1/T_2$ Multi-Run Drift and Control-Arm Arbitration}}

\author{
  \textbf{M\'at\'e R\'obert}\thanks{Corresponding author: \texttt{materobi86@gmail.com}. Open to research collaborations and academic partnerships.}\\
  \textit{Independent Quantum Researcher, Silent Coder Lab}\\
  \textit{Budapest, Hungary}
}
\date{November 2026}

\begin{document}
\maketitle

\begin{abstract}
Daily backend calibration reports on cloud-hosted superconducting transmon processors~\cite{krantz_guide,blais_cqed} provide reference decoherence lifetimes ($T_1$ energy relaxation and Hahn-echo $T_2$). However, quantum memory circuits and state-preservation subroutines frequently execute general unitary operations with arbitrary rotation axes rather than idealized, resonant Clifford refocussing pulses~\cite{hahn1950,carr_purcell1954}. In this study, we benchmark the decoherence of an arbitrary-axis unitary inversion sequence ($U \cdot \text{delay}(\tau) \cdot U^{\dagger}$) across delay intervals $\tau \in [0, 800]\,\mu\text{s}$ on physical qubit 91 of the 156-qubit IBM Heron processor (\texttt{ibm\_kingston})~\cite{heron2024}. Crucially, we incorporate an independent $T_1$ control arm ($U \cdot \text{delay}(\tau)$ without inversion) as an in-situ clock reference to verify hardware timescale fidelity. We show that while the average measured $T_1$ ($274.0 \pm 36.6\,\mu\text{s}$ across multi-run sweeps) verifies the compiler's pulse duration scaling, physical lifetimes undergo pronounced inter-run variations ($T_1 \in [211.6, 309.3]\,\mu\text{s}$, a $46\%$ excursion). Across replicate sweeps, effective random-axis echo lifetimes drop to $T_{2,\text{eff}} = 89.6 \pm 39.2\,\mu\text{s}$ (down to $43.1\,\mu\text{s}$ in degraded runs), substantially lower than the quoted daily calibration baseline ($T_{2,\text{cal}} = 121.3\,\mu\text{s}$). Replicate sweeps separated by 75-second idling windows demonstrate severe excess variance ($\chi^2 = 60.9$, with inter-run variance $3.6\times$ Poisson shot noise, $p \approx 4.5 \times 10^{-12}$). Rather than isolated dephasing, our multi-run analysis reveals correlated $T_1/T_2$ drift consistent with two-level system (TLS) spectral diffusion~\cite{burnett2019,muller2019tls}. We conclude that calibration-quoted $T_2$ figures systematically overestimate operational state retention in arbitrary-axis quantum algorithms.
\end{abstract}

\vspace{1em}
\noindent \textbf{Keywords:} Superconducting transmons, Decoherence benchmarking, Hahn echo, Spectral diffusion, Two-level systems, Quantum memory, Hardware characterization.

\section{Introduction}
Superconducting circuit architectures based on transmon qubits~\cite{koch2007transmon} have emerged as leading platforms for programmable quantum computation~\cite{arute2019supremacy,blais_cqed}. IBM Quantum and other cloud providers publish automated daily calibration snapshots reporting energy relaxation times ($T_1$) and spin-echo coherence times ($T_2$)~\cite{krantz_guide}. These metrics are traditionally acquired via specialized calibration schedules using symmetric resonant $\pi$-pulses aligned strictly along orthogonal Pauli axes ($X$ or $Y$)~\cite{hahn1950,meiboom1958}.

In practical quantum memory registers and algorithmic subroutines, arbitrary state vectors $|\psi\rangle = U|0\rangle$ undergo variable delay intervals before algorithmic inversion via $U^{\dagger}$~\cite{nielsen_chuang,viola1998dd}. A frequent operational heuristic assumes that such reversal operations act as dynamical decoupling echoes. In this work, we test this assumption through systematic multi-run hardware experiments on IBM Heron (\texttt{ibm\_kingston}) using physical qubit 91.

\section{Experimental Architecture: The $T_1$ Arbitration Control Arm}
Decoherence measurements were conducted on physical qubit 91 of the 156-qubit Heron r3 processor \texttt{ibm\_kingston}~\cite{heron2024}. For each delay duration $\tau \in \{0, 50, 100, 200, 300, 400, 600, 800\}\,\mu\text{s}$, two parallel circuit branches were executed with 1024 shots per setting:

\begin{enumerate}
  \item \textbf{Unitary Inversion Branch (Echo):} State preparation $U$, followed by idling delay $\tau$, followed by inverse unitary $U^{\dagger}$ and computational basis readout. In the ideal noiseless limit, outcome probability evaluates strictly to $P(|0\rangle) = 1.0$.
  \item \textbf{Decay Control Branch (Arbitration Arm):} State preparation $U$, followed by idling delay $\tau$, and direct readout without inversion~\cite{bylander2011noise}.
\end{enumerate}

The control arm decays purely under longitudinal relaxation $T_1$. Because $T_1$ decay is insensitive to pure phase dephasing ($T_\phi$), it serves as an empirical arbitration standard: if the compiler fails to properly scale the hardware delay parameter ($\text{dt}$), the extracted $T_1$ diverges from the physical baseline.

\begin{table}[htbp]
\centering
\begin{tabular}{lcccc}
\toprule
\textbf{Sweep Run} & \textbf{$T_{2,\text{eff}}$ ($\mu\text{s}$)} & \textbf{$T_{1,\text{eff}}$ ($\mu\text{s}$)} & \textbf{Ratio $T_2/T_1$} & \textbf{Drift State} \\
\midrule
Run 1 & $77.5 \pm 12.9$ & $285.5 \pm 27.6$ & $0.27$ & Nominal \\
Run 2 & $43.1 \pm 9.0$  & $211.6 \pm 32.0$ & $0.20$ & Severe TLS dip \\
Run 3 & $140.8 \pm 10.7$& $272.6 \pm 24.0$ & $0.52$ & High coherence \\
Run 4 & $117.7 \pm 12.5$& $309.3 \pm 24.0$ & $0.38$ & Recovered \\
Run 5 & $69.1 \pm 9.0$  & $290.3 \pm 13.2$ & $0.24$ & Marginal \\
\midrule
\textbf{Multi-run Mean} & $\mathbf{89.6 \pm 39.2}$ & $\mathbf{274.0 \pm 36.6}$ & $0.32$ & $\chi^2 = 60.9$ ($p \approx 4.5\times 10^{-12}$) \\
Calibration Reference & $121.3$ & $132.4$ & $0.92$ & Backend snapshot \\
\bottomrule
\end{tabular}
\caption{Multi-run decoherence metrics on \texttt{ibm\_kingston} (Qubit 91). Runs 1--5 demonstrate correlated $T_1/T_2$ multi-run drift, rejecting the hypothesis of stationary dephasing.}
\label{tab:qds_multirun}
\end{table}

\section{Results and Correlated Drift Analysis}
As documented in Table~\ref{tab:qds_multirun}, the control arm verifies the microsecond time base of the pulse generator. However, two critical physical phenomena emerge from the multi-run dataset:

\subsection{Systematic Echo Deficit}
The measured mean effective dephasing time across five runs is $T_{2,\text{eff}} = 89.6 \pm 39.2\,\mu\text{s}$, compared to the daily calibration report of $T_{2,\text{cal}} = 121.3\,\mu\text{s}$~\cite{krantz_guide}. Arbitrary-axis inversion circuits $U \cdot \text{delay} \cdot U^\dagger$ cannot form a closed refocussing trajectory on the Bloch sphere against non-commuting background noise Hamiltonians~\cite{viola1998dd}, leaving residual low-frequency dephasing uncompensated.

\subsection{Correlated $T_1/T_2$ Drift and TLS Fluctuations}
Prior analyses of superconducting qubit stability often treat $T_1$ relaxation as stationary while attributing phase fluctuations solely to flux noise~\cite{koch2007transmon}. In our measurements, $T_1$ varies over a $46\%$ window ($211.6\,\mu\text{s}$ to $309.3\,\mu\text{s}$). Notably, Run 2 manifests a simultaneous minimum in both $T_1$ ($211.6\,\mu\text{s}$) and $T_2$ ($43.1\,\mu\text{s}$). This co-variation demonstrates that the qubit is coupled to fluctuating microscopic two-level systems (TLS) residing in the dielectric substrate or Josephson junction interfaces~\cite{burnett2019,muller2019tls,place2021}. Spectral diffusion of near-resonant defect modes modulates both the purcell/relaxation decay channels and pure dephasing rates simultaneously.

Across 4 replicate sweeps separated by 75-second idle intervals, the measured survival probabilities yield $\chi^2 = 60.9$ (degrees of freedom $= 12$, $p \approx 4.5 \times 10^{-12}$), showing inter-run variance $3.6\times$ beyond the Poisson binomial shot noise floor. This demonstrates that calibration-table snapshots represent transient operating conditions rather than stationary guarantees.

\section{Conclusion}
Calibration-quoted $T_2$ values obtained via standard Hahn echo schedules cannot be assumed for general unitary algorithms. Quasi-static $1/f$ drift and fluctuating TLS defects cause correlated $T_1/T_2$ excursions on sub-minute timescales. Algorithmic error budgets and quantum memory designs must incorporate multi-run dynamical arbitration arms rather than static calibration tables.

\section*{Data and Code Availability}
All raw hardware JSON execution logs, pulse level definitions, and offline validation suites (38 unit and statistical tests) are publicly available under the MIT license at: \url{https://github.com/silentnoisehun/quantum-delayed-sync}.

\begin{thebibliography}{99}
\bibitem{blais_cqed}
A.~Blais, A.~L.~Grimsmo, S.~M.~Girvin, and A.~Wallraff,
``Circuit quantum electrodynamics,''
\emph{Reviews of Modern Physics}, vol.~93, no.~2, p.~025005, 2021.

\bibitem{krantz_guide}
P.~Krantz, M.~Kjaergaard, F.~Yan, T.~P.~Orlando, S.~Gustavsson, and W.~D.~Oliver,
``A quantum engineer's guide to superconducting qubits,''
\emph{Applied Physics Reviews}, vol.~6, no.~2, p.~021318, 2019.

\bibitem{koch2007transmon}
J.~Koch, T.~M.~Yu, J.~Gambetta, A.~A.~Houck, D.~I.~Schuster, J.~Majer, A.~Blais, M.~H.~Devoret, S.~M.~Girvin, and R.~J.~Schoelkopf,
``Charge-insensitive qubit design in the transmon regime,''
\emph{Physical Review A}, vol.~76, no.~4, p.~042319, 2007.

\bibitem{hahn1950}
E.~L.~Hahn,
``Spin echoes,''
\emph{Physical Review}, vol.~80, no.~4, pp.~580--594, 1950.

\bibitem{carr_purcell1954}
H.~Y.~Carr and E.~M.~Purcell,
``Effects of diffusion on free precession in nuclear magnetic resonance experiments,''
\emph{Physical Review}, vol.~94, no.~3, pp.~630--638, 1954.

\bibitem{meiboom1958}
S.~Meiboom and D.~Gill,
``Modified spin-echo method for measuring nuclear relaxation times,''
\emph{Review of Scientific Instruments}, vol.~29, no.~8, pp.~688--691, 1958.

\bibitem{viola1998dd}
L.~Viola and S.~Lloyd,
``Dynamical suppression of decoherence in two-state quantum systems,''
\emph{Physical Review A}, vol.~58, no.~4, pp.~2733--2744, 1998.

\bibitem{bylander2011noise}
J.~Bylander, S.~Gustavsson, F.~Yan, F.~Yoshihara, K.~Harrabi, G.~Fitch, D.~G.~Cory, Y.~Nakamura, J.-S.~Tsai, and W.~D.~Oliver,
``Noise spectroscopy through dynamical decoupling with a superconducting flux qubit,''
\emph{Nature Physics}, vol.~7, no.~7, pp.~565--570, 2011.

\bibitem{burnett2019}
J.~J.~Burnett, A.~Bengtsson, M.~Scigliuzzo, D.~Niepce, M.~Kudra, P.~Delsing, and J.~Bylander,
``Decoherence benchmarking of superconducting qubits,''
\emph{npj Quantum Information}, vol.~5, no.~1, p.~54, 2019.

\bibitem{muller2019tls}
C.~M\"uller, J.~H.~Cole, and J.~Lisenfeld,
``Towards understanding two-level-systems in amorphous solids: insights from quantum circuits,''
\emph{Reports on Progress in Physics}, vol.~82, no.~12, p.~124501, 2019.

\bibitem{place2021}
A.~P.~M.~Place et~al.,
``New material platform for superconducting transmon qubits with coherence times exceeding 0.3 milliseconds,''
\emph{Nature Communications}, vol.~12, no.~1, p.~1779, 2021.

\bibitem{arute2019supremacy}
F.~Arute et~al.,
``Quantum supremacy using a programmable superconducting processor,''
\emph{Nature}, vol.~574, no.~7779, pp.~505--510, 2019.

\bibitem{nielsen_chuang}
M.~A.~Nielsen and I.~L.~Chuang,
\emph{Quantum Computation and Quantum Information},
Cambridge University Press, 2010.

\bibitem{heron2024}
IBM Quantum,
``IBM Quantum Heron processor architecture and calibration specifications,''
IBM Technical Whitepaper, 2024.
\end{thebibliography}

\end{document}
